\chapter{Implementation Details}
\label{ch:implementation}

This section translates the algorithms of Chapter~\ref{ch:methodology} into concrete artefacts: the proxy server, the typed TypeScript SDK, the backend service, and the Cloud Run deployment.

\section{Repository Layout}
\label{sec:repo}

The platform is organised as a hybrid monorepo:
\begin{itemize}
    \item \texttt{sprout-web/} --- Next.js 14 frontend (TypeScript, Tailwind, Shadcn/ui, Zustand). Static export deployed to Firebase Hosting.
    \item \texttt{sprout-backend/} --- Firebase Cloud Functions (Node.js + TypeScript). REST API, Firestore triggers, auth middlewares.
    \item \texttt{gemini-proxy/} --- Stateless Node.js WebSocket gateway on Cloud Run. No per-user data beyond the in-process session map.
    \item \texttt{packages/gemini-proxy-sdk/} --- Published as \texttt{@sprout/gemini-proxy-sdk}.
    \item \texttt{sprout-app/} --- Terraform provisioning Firebase, Cloud Run, Cloud Tasks, and secrets~\cite{firebase2024,cloudrun2024}.
\end{itemize}

\section{WebSocket Wire Protocol}
\label{sec:wss-protocol}

Clients connect with a signed JWT in the query string because browser \code{WebSocket} cannot attach arbitrary headers. Listing~\ref{lst:wss-connect} shows the URL pattern; Listing~\ref{lst:wss-frame} shows the JSON envelope for non-audio frames.

\begin{lstlisting}[caption={WebSocket upgrade URL (browser).},label=lst:wss-connect]
const url = new URL(process.env.NEXT_PUBLIC_GEMINI_PROXY_URL);
url.searchParams.set('auth', await tokenProvider.getToken());
const socket = new WebSocket(url.toString());
\end{lstlisting}

\begin{lstlisting}[caption={Opaque JSON control frame (client $\rightarrow$ proxy).},label=lst:wss-frame]
{
  "type": "client_content",
  "modality": "text",
  "payload": { "text": "Explain async/await in JavaScript." }
}
\end{lstlisting}

Audio and video use binary frames with a fixed 20\,ms PCM chunk header prepended by the SDK; the proxy forwards them without parsing content (Algorithm~\ref{alg:a2}). Table~\ref{tab:wss-frames} summarises frame categories.

\tablesetup
\begin{table}[H]
\centering
\caption{WebSocket frame categories handled by the proxy.}
\label{tab:wss-frames}
\small
\begin{tabular}{|p{2.4cm}|p{1.5cm}|p{8.4cm}|}
\hline
\chead{Direction} & \chead{Type} & \chead{Handling} \\
\hline
Client $\rightarrow$ proxy & Binary & Forward to Gemini Live; update byte counters \\
\hline
Client $\rightarrow$ proxy & JSON & Route text, tool ack, heartbeat \\
\hline
Proxy $\rightarrow$ client & Binary & Model audio/stream chunks \\
\hline
Proxy $\rightarrow$ client & JSON & Transcript, \code{usageMetadata}, \code{function\_call} \\
\hline
Either & Close & Map to Table~\ref{tab:close-codes}; persist \code{closeCode} \\
\hline
\end{tabular}
\end{table}

\section{Gemini Proxy Server}
\label{sec:proxy}

The proxy is built around the \texttt{ws} WebSocket library. Listing~\ref{lst:jwt} shows JWT validation; Listing~\ref{lst:session} the session manager; Listing~\ref{lst:backend} the backend event client.

\begin{lstlisting}[caption={JWT validation middleware (\texttt{auth-middleware.js}).},label=lst:jwt]
const { OAuth2Client } = require('google-auth-library');
const client = new OAuth2Client(process.env.GOOGLE_CLIENT_ID);

async function validateGoogleJWT(token) {
  const ticket = await client.verifyIdToken({
    idToken: token,
    audience: process.env.GOOGLE_CLIENT_ID,
  });
  return ticket.getPayload();
}
\end{lstlisting}

\begin{lstlisting}[caption={Session manager (\texttt{session-manager.js}).},label=lst:session]
class SessionManager {
  constructor() {
    this.sessions = new Map();
    this.userSessions = new Map();
  }
  createSession(uid, userData) {
    const sessionId = generateSessionId();
    const session = {
      id: sessionId, uid, userData,
      startTime: Date.now(),
      tokenUsage: { input: 0, output: 0 },
      messageCount: 0, status: 'active'
    };
    this.sessions.set(sessionId, session);
    this.userSessions.set(uid, sessionId);
    return session;
  }
}
\end{lstlisting}

\begin{lstlisting}[caption={Backend event client.},label=lst:backend]
class BackendClient {
  async sendSessionEvent(sessionId, eventType, data) {
    try {
      await axios.post(
        `${this.baseUrl}/api/sessions/${sessionId}/events`,
        { type: eventType, data, timestamp: Date.now() }
      );
    } catch (e) {
      console.error('Failed to send backend event:', e);
    }
  }
}
\end{lstlisting}

\section{Client SDK}
\label{sec:sdk}

The frontend uses \texttt{@sprout/gemini-proxy-sdk}, which standardises authentication, reconnection, audio framing, and metric capture. Listing~\ref{lst:sdk-api} shows the public TypeScript surface; Listing~\ref{lst:browser} typical browser usage.

\begin{lstlisting}[caption={Public SDK API.},label=lst:sdk-api]
export interface TokenProvider {
  getToken(): Promise<string>;
}
export interface GeminiProxyClientOptions {
  proxyUrl: string;
  tokenProvider: TokenProvider;
  reconnect?: { enabled: boolean; maxRetries?: number };
}
export interface GeminiProxyClient {
  connect(): Promise<void>;
  disconnect(code?: number, reason?: string): void;
  send(data: unknown): void;
  sendText(text: string): void;
  on(listener: (e: GeminiEvent) => void): () => void;
}
\end{lstlisting}

\begin{lstlisting}[caption={Browser usage.},label=lst:browser]
const client = createGeminiClient({
  proxyUrl: process.env.NEXT_PUBLIC_GEMINI_PROXY_URL!,
  tokenProvider: { getToken: () => auth.currentUser!.getIdToken() },
  reconnect: { enabled: true, maxRetries: 5, backoffMs: 1000 },
});
await client.connect();
client.on((e) => {
  if (e.type === 'message') { /* update UI */ }
  if (e.type === 'usage') { /* update token meter */ }
});
\end{lstlisting}

\section{Backend Service Surface}
\label{sec:backend-api}

Public endpoints include \texttt{POST /onboarding/skills}, \texttt{POST /courses/knowledge-map}, \texttt{GET /courses/\{id\}/lessons/\{id\}/completion}, and \texttt{POST /conversations/feedback}. Internal proxy endpoints include \texttt{POST /internal/session/create-session} and \texttt{POST /internal/chat/events/callback}. The split ensures the live data plane never invokes user-scoped auth paths.

Table~\ref{tab:api-catalogue} catalogues the control-plane surface area most relevant to the AI subsystem. All public routes require a Firebase ID token unless noted.

\tablesetup
\begin{table}[H]
\centering
\caption{Representative REST API catalogue (\texttt{sprout-backend}).}
\label{tab:api-catalogue}
\small
\begin{tabular}{|p{4.6cm}|p{1.2cm}|p{7.5cm}|}
\hline
\chead{Route} & \chead{Verb} & \chead{Purpose} \\
\hline
\texttt{/onboarding/skills} & POST & NL skill discovery; LLM-refined suggestions \\
\texttt{/courses/knowledge-map} & POST & Enqueue knowledge-map worker (Cloud Tasks) \\
\texttt{/courses/\{id\}/lessons} & GET & List lessons and completion flags \\
\texttt{/courses/\{id\}/lessons/\{lid\}/completion} & GET & Poll lesson completion (Algorithm~\ref{alg:a4}) \\
\texttt{/conversations/feedback} & POST & Thumbs up/down; triggers Algorithm~\ref{alg:a5} \\
\texttt{/internal/session/create-session} & POST & Proxy-only session metadata registration \\
\texttt{/internal/chat/events/callback} & POST & Function-call and completion callbacks from proxy \\
\texttt{/admin/users/\{uid\}/ban} & POST & Set \texttt{users.banned}; enforced at proxy (A1) \\
\hline
\end{tabular}
\end{table}

Table~\ref{tab:gemini-session-fields} documents the \texttt{gemini\_sessions} document written by the proxy and backend aggregators. These fields power billing dashboards and operational queries.

\tablesetup
\begin{table}[H]
\centering
\caption{\texttt{gemini\_sessions} field reference.}
\label{tab:gemini-session-fields}
\small
\begin{tabular}{|p{3.0cm}|p{1.9cm}|p{8.3cm}|}
\hline
\chead{Field} & \chead{Type} & \chead{Semantics} \\
\hline
\texttt{uid} & string & Firebase user or anonymous UID \\
\texttt{startTime} & timestamp & WebSocket accept time (proxy clock) \\
\texttt{endTime} & timestamp & Close or idle timeout \\
\texttt{tokenUsage.input} & int & Cumulative prompt tokens from \texttt{usageMetadata} \\
\texttt{tokenUsage.output} & int & Cumulative response tokens \\
\texttt{messageCount} & int & Frames routed (Algorithm~\ref{alg:a2}) \\
\texttt{modality} & string & \texttt{voice}, \texttt{text}, \texttt{webcam}, \texttt{screen} \\
\texttt{status} & string & \texttt{active}, \texttt{closed}, \texttt{error} \\
\texttt{closeCode} & int & WebSocket close code (Table~\ref{tab:close-codes}) \\
\hline
\end{tabular}
\end{table}

\section{\texttt{course\_profiles} (Planned)}
\label{sec:course-profiles-impl}

The deployed backend today writes a single \texttt{conversations/\{id\}.summary} per feedback event (Algorithm~\ref{alg:a5}). Section~\ref{sec:multi-intent} extends this with a per-course subcollection so session priming can select top-$k$ intent profiles instead of one collapsed persona. The change is confined to Firestore triggers and prompt assembly; the Gemini Proxy transport is unchanged.

Listing~\ref{lst:course-profile-doc} shows the planned document shape under \texttt{users/\{uid\}/course\_profiles/\{courseId\}}. The feedback trigger will upsert this document alongside the conversation summary.

\begin{lstlisting}[caption={Planned Firestore document for \texttt{course\_profiles/\{courseId\}} (JSON).},label=lst:course-profile-doc]
// users/{uid}/course_profiles/{courseId}
{
  "rollingSummary": "Weak on async/await; strong on DOM APIs.",
  "topics": ["async/await", "closures", "DOM"],
  "masteryScore": 0.62,
  "lastActiveAt": "2026-05-19T08:30:00.000Z",
  "embedding": null
}
\end{lstlisting}

Listing~\ref{lst:build-context} sketches the backend call site before opening a lesson session: \textsc{BuildSessionContext} (Algorithm~\ref{alg:context}) merges the active course profile with up to $k-1$ recent alternatives, then injects the result into the Gemini system prompt.

\begin{lstlisting}[caption={Planned session priming hook (\texttt{sprout-backend}).},label=lst:build-context]
async function primeLessonSession(uid: string, courseId: string, lessonId: string) {
  const ctx = await buildSessionContext(uid, courseId, lessonId, { k: 2 });
  return {
    systemInstruction: formatSystemPrompt(ctx),
    lessonId,
    courseId,
  };
}
\end{lstlisting}

Rollout is incremental: ship schema and trigger writes first (dual-write to \texttt{summary} and \texttt{course\_profiles}), then switch priming to \texttt{buildSessionContext} behind a feature flag for the Phase~3 A/B in Section~\ref{sec:online-protocol}.

\section{Cloud Run Deployment}
\label{sec:cloudrun}

Listing~\ref{lst:cloudrun} shows the service definition: container concurrency 50, CPU throttling disabled, minimum instances 1 to avoid cold-start latency.

\begin{lstlisting}[caption={Cloud Run service definition for \texttt{gemini-proxy}.},label=lst:cloudrun]
apiVersion: serving.knative.dev/v1
kind: Service
metadata:
  name: gemini-proxy
spec:
  template:
    metadata:
      annotations:
        run.googleapis.com/cpu-throttling: 'false'
        run.googleapis.com/min-instances: '1'
        run.googleapis.com/max-instances: '200'
    spec:
      containerConcurrency: 50
      containers:
      - image: gcr.io/$PROJECT_ID/gemini-proxy:latest
        ports:
        - name: h2c
          containerPort: 8080
\end{lstlisting}

\section{Configuration and Secrets}
\label{sec:secrets}

Credentials are sourced from Google Secret Manager. Non-sensitive configuration flows from Terraform variables; \texttt{NEXT\_PUBLIC\_*} variables are injected at build time and are not secrets.

\section{Frontend Components}
\label{sec:frontend}

The AI surface is concentrated in \texttt{AssistantConnector} (SDK lifecycle and modality selection), \texttt{AssistantChat} (transcript and text input), and \texttt{SkillMapChart} (React Flow knowledge map). The \texttt{useLessonCompletionPolling} hook implements Algorithm~\ref{alg:a4} on the Firestore client SDK.
